\section{Introduction}

A household robot that learns to wipe a counter from one recording will wipe that counter. Change the laminate from white to blue, swap the bucket, or move the task to a different kitchen, and a policy trained on the single clip has no reason to generalize. Vision-language-action (VLA) models such as RT-2~\cite{rt2}, OpenVLA~\cite{openvla} and $\pi_0$~\cite{pi0} get around this by training on very large demonstration sets~\cite{oxe}, and a large share of that scale is spent on appearance rather than on new skills. Collecting appearance diversity by hand is the slow part. Our source material for this work is 8 egocentric recordings of ordinary household chores, about 80 minutes of raw footage, and even that small set took deliberate effort to stage and record.

The two existing routes to cheap appearance diversity each give up something. Domain randomization~\cite{tobin2017,peng2018} perturbs colors, textures and lighting at random. It is free and it works for sim-to-real transfer, but it has no notion of the task: it will recolor a cleaning bucket into a texture no bucket has, and it treats a rag on a counter the same as a rag on a shelf. Semantic augmentation with image diffusion models~\cite{rosie,genaug} fixes the plausibility problem for single frames, but robot learning increasingly consumes video, and per-frame editing does not hold objects and hands steady across time. Video-to-video models conditioned on structure~\cite{gen1,controlnet} can restyle a whole clip, and NVIDIA's Cosmos family~\cite{cosmos2025,cosmostransfer1} packages this for physical AI. What none of these supply is a decision procedure: given one clip, which objects should vary, in which ways, and how does one keep everything else fixed.

We treat that decision as a reasoning problem and the rendering as a constrained generation problem, and we chain existing models to solve each. A reasoning world foundation model reads the clip and writes a structured ontology: the action primitive, the task-relevant objects with bounding boxes, each object's current appearance and state, and the variations that would matter for a robot. A language model filters those variations for physical plausibility and turns the survivor into a generation prompt. Tracked segmentation masks, built by subtracting the objects on a surface from the surface itself, feed a multi-control video-to-video model together with edge, depth and blurred-RGB maps. The key idea is that the ontology both chooses the variation and names what must not change, so the same structured output drives the prompt and the mask.

The paper reports the following.
\begin{itemize}
    \item A pipeline (Fig.~\ref{fig:pipeline}) that takes a 15-second egocentric clip to a rendered appearance variant with no manual step other than picking which variation to render, with all prompts and the ontology schema released.
    \item A measured control-weight configuration for Cosmos Transfer~2.5 on egocentric manipulation video, and the failure mode that pins each weight: lowering the edge weight below 0.5 collapses object shape, raising depth above 0.2 flattens the scene, and the blurred-RGB control must sit near zero or it reverts the recolor.
    \item A quantitative finding that soft multi-control conditioning does not confine edits. On our best case the target region changes by a mean $\Delta E$ of 75 while the protected region changes by 35, and 95\% of protected pixels change visibly.
    \item A negative result on action-level variation: Cosmos Predict2 given a manipulation clip and an action prompt produced no visible change in the manipulated object.
\end{itemize}

The reasoning stage costs 130~s per clip on one A100 and yields 72 admissible variations on average; one Transfer render costs 20~min. We rendered 64 variants. At those rates the 116 clips we hold could become about 8{,}400, but the leakage numbers above say that the variants would not yet be the controlled edits a robot-learning curriculum needs, and the rest of the paper is about why.
